% Supplementary Information for MDArena.
% Adapted from the supplied SI: separate title page, single-column text,
% blue links, S-prefixed pages, figures, tables, and equations.
% Standard article class replaces the journal-specific ACS class.
% Content adapted from the supplied NeurIPS workshop appendix.
% Title and author metadata follow the current Digital Discovery manuscript.
% Workshop task descriptions and comparisons are preserved, not re-audited.
% REVIEW: Task 39 describes explicit ligand B-state perturbations as required;
% check against the evaluated task version before journal submission.
% REVIEW: Confirm the Task 22 supporting-plot verification description.
% REVIEW: Update comparison-table classifications against the cited sources.
\documentclass[11pt,a4paper]{article}
\usepackage[T1]{fontenc}
\usepackage[utf8]{inputenc}
\usepackage[margin=2.5cm]{geometry}
\usepackage{graphicx}
\usepackage{amsmath,amssymb}
\usepackage{booktabs,longtable,array}
\usepackage{float}
\usepackage{xcolor}
\usepackage[numbers,sort&compress]{natbib}
\newcolumntype{L}[1]{>{\raggedright\arraybackslash}p{#1}}
\newcommand{\cmark}{\ensuremath{\checkmark}}
\newcommand{\xmark}{\ensuremath{\times}}
\newcommand{\pmark}{\ensuremath{\circ}}
\usepackage{hyperref}
\usepackage[capitalize,noabbrev]{cleveref}
% Missing external assets are marked visibly until the original files are added.
\newcommand{\suppfigure}[1]{%
  \IfFileExists{#1}{\includegraphics[width=\textwidth]{#1}}{%
    \fbox{\parbox[c][4cm][c]{0.93\textwidth}{\centering
      Figure file not supplied.\\[1ex]\texttt{\detokenize{#1}}}}}}

\hypersetup{colorlinks=true,linkcolor=blue,citecolor=blue,urlcolor=blue,
 pdftitle={Supplementary Information: Assessing coding-agent reliability in molecular dynamics workflows with MDArena}}
\renewcommand{\thepage}{S\arabic{page}}
\renewcommand{\thefigure}{S\arabic{figure}}
\renewcommand{\thetable}{S\arabic{table}}
\renewcommand{\theequation}{S\arabic{equation}}
\newcommand{\draftnote}[1]{\par\medskip\noindent\textcolor{red!65!black}{\textbf{Draft note:} #1}\par\medskip}
% Optional references to main.tex: uncomment when main.aux is available.
% \usepackage{xr}
% \externaldocument[main-]{main}
\begin{document}
\begin{titlepage}
\centering
\vspace*{1cm}
{\Large\bfseries Supplementary Information:\par}
\vspace{0.4cm}
{\LARGE\bfseries Assessing coding-agent reliability in molecular dynamics workflows with MDArena\par}
\vspace{1cm}
{\large Nithishwer Mouroug Anand\textsuperscript{a,*,$\ddagger$},
Wei-Tse Hsu\textsuperscript{a,$\ddagger$},
Kyle Vaccaro\textsuperscript{a,b},
Eden James Gage\textsuperscript{a},
Jonathan David Colburn\textsuperscript{a},
Linda Xi Phan\textsuperscript{a},
Minjoon Seo\textsuperscript{a},
Kevin Guan\textsuperscript{a}, and
Philip C. Biggin\textsuperscript{a,*}\par}
\vspace{0.8cm}
\textsuperscript{a}University of Oxford\par
\vspace{0.2cm}
\textsuperscript{b}Scripps Research Institute\par
\vspace{0.8cm}
\textsuperscript{$\ddagger$}Nithishwer Mouroug Anand and Wei-Tse Hsu contributed equally to this work.\par
\vspace{0.5cm}
\textsuperscript{*}Corresponding authors:\par
\href{mailto:nithishwer.mourouganand@reuben.ox.ac.uk}{nithishwer.mourouganand@reuben.ox.ac.uk}\par
\href{mailto:philip.biggin@bioch.ox.ac.uk}{philip.biggin@bioch.ox.ac.uk}\par
% Affiliations intentionally reproduce the current manuscript; expand both
% documents together if departmental names and postal addresses are required.
\vfill
\end{titlepage}
\setcounter{page}{2}
\tableofcontents
\clearpage



\section{Representative Tasks}
\label{app:representative_tasks}

To illustrate the nature of MDArena beyond its aggregate composition, we present three representative tasks exhibiting distinct forms of difficulty: extending a force field, constructing a charged-ligand free-energy workflow, and evaluating the suitability of a scientific method. The excerpts below omit only standardized file paths, output specifications, and time limits; complete prompts are provided with the benchmark.

\subsection{Parameterizing a non-standard phosphoaspartate residue}
\label{app:task47}

\paragraph{Task.}
The agent is given a pentapeptide from a P-type ATPase phosphoenzyme intermediate containing a dianionic phosphoaspartate residue, together with an unmodified CHARMM36 force field and literature parameters for the modified side chain:

\begin{quote}
Extend a copy of the CHARMM36 force field so that it recognizes the D2P residue and can build a complete, chemically sensible topology with the correct connectivity, atom types, charges, and bonded parameters. Only retype or recharge the atoms and add the bonded terms specified by the supplied literature parameters; all other terms should be retained from the standard aspartate residue. Generate a topology for the peptide, solvate the system, and perform energy minimization followed by NVT equilibration.
\end{quote}

\paragraph{Why it is challenging.}
Unlike more routinely supported phosphorylated residues such as phosphoserine and phosphotyrosine, phosphoaspartate is not included in the standard CHARMM36 residue set supplied with GROMACS. The task therefore provides the dianionic phosphoaspartate parameters reported by Damjanović et al. for simulations of phosphorylated NtrC~\citep{damjanovic2009self}. The agent must infer the esterified oxygen from the molecular structure to correctly assign atom types and charges, preserve unrelated parameters of standard aspartate, and convert CHARMM units and functional forms into their GROMACS equivalents. Importantly, the prompt provides the target literature parameters but does not prescribe how to resolve these mapping and conversion issues. The agent must also diagnose the residue-classification issue that prevents \texttt{pdb2gmx} from processing the peptide, construct a self-contained force-field extension, and carry the resulting system through topology generation, solvation, minimization, and equilibration. Despite being conceptually straightforward, this task is operationally error-prone even for human researchers, and is therefore particularly interesting: an agent capable of performing such work reliably would provide practical value.

\paragraph{Verification.}
The deterministic verifier comprehensively examines the D2P connectivity, atom types, charges, and bonded parameters. It confirms the validity of the agent's submitted force field by reconstructing the complete preparation stages, including minimization and equilibration, in a fresh directory. In addition, it checks that the other parameters in the CHARMM36 force field are unchanged, and whether submitted structures are consistent with completed minimization and NVT runs to prevent falsification. The process verifier evaluates how the agent resolved the residue connectivity and classification, converted the literature parameters, preserved the supplied inputs, and substantiated its claim that each stage completed successfully.

\subsection{Setting up a binding free energy calculation for a charged D3 receptor ligand}
\label{app:task39}

\paragraph{Task.}

The agent is given prepared GROMACS inputs for the human dopamine D3 receptor membrane system with a (+1)-charged ligand 11a:

\begin{quote}
Build the complex leg of a single-ligand absolute binding free-energy setup using the prescribed restraint, Coulomb, and van der Waals lambda schedules. Use \texttt{mdrestraintsgenerator} to generate Boresch restraints and select the closest restraint frame for the complex setup. Generate all required simulation stages and compile the energy-minimization input for every window, but do not run dynamics or replica-exchange simulations.
\end{quote}

\paragraph{Why it is challenging.}
Although the prompt specifies the lambda schedules and output structure in detail, it deliberately does not state how charge neutrality should be maintained. Annihilating ligand 11a changes the net charge of the system, so the agent must recognize the need for a compensating alchemical counterion, encode its neutralization across the Coulomb and van der Waals phases, and include it consistently in the required index groups. The ligand must likewise be perturbed explicitly through B-state charges and dummy atom types rather than solely through \texttt{couple-moltype}, which cannot represent the intended coupled ligand--ion transformation. The task therefore tests whether an agent can identify an unstated but essential physical requirement while implementing an otherwise detailed specification.

\paragraph{Verification.}
The deterministic verifier carries out four substantive checks. Specifically, it (1) examines the generated Boresch restraints and closest restraint frame, (2) checks the complete window structure and lambda schedules, (3) confirms the explicit ligand and counterion perturbations and corresponding index groups, and (4) tests whether representative windows can be processed by \texttt{grompp} without unexpected warnings. The process verifier assesses whether the agent used the supplied trajectory to generate the restraints, recognized the need for a compensating counterion, implemented the perturbations explicitly, and built the workflow from the supplied inputs.

\subsection{Assessing a PCA-derived collective variable}
\label{app:task22}

\paragraph{Task.}
The agent is given an unbiased trajectory of the PepT2 membrane transporter and PLUMED definitions of two structural variables describing its gate opening and closing:

\begin{quote}
Analyze whether the dominant C-alpha PCA motion in the PepT2 trajectory is a suitable collective variable for enhanced sampling of the PepT2 opening/closing motion. Report the PCA setup, including alignment and atom selection; the evidence used to decide whether PCA captures opening and closing; at least one supporting plot; and a binary recommendation to proceed with the PCA-derived CV or try a different CV.
\end{quote}

\paragraph{Why it is challenging.}
The dominant principal component captures the largest structural variance, but it is not necessarily a useful reaction coordinate for the conformational transition of interest. The agent must therefore align the trajectory appropriately, perform principal component analysis (PCA) on C-alpha atoms, independently calculate the supplied opening and closing variables, compare their distributions and correlation, and interpret whether the PCA coordinate captures the relevant motion. Success requires a scientifically justified decision rather than only producing a valid trajectory analysis.

\paragraph{Verification.}
Deterministic checks require the correct recommendation and a supporting plot referenced in the final answer. Complementary process criteria assess whether the agent actually performed the PCA, evaluated the supplied structural variables, quantified their relationship, and based its recommendation on this evidence. This prevents an agent from receiving full credit by guessing the correct conclusion without conducting the required analysis.

\section{Supplementary Tables}

\begingroup
\scriptsize
\setlength{\tabcolsep}{2pt}
\renewcommand{\arraystretch}{1.12}
\begin{longtable}{@{}L{0.19\textwidth}L{0.15\textwidth}L{0.29\textwidth}L{0.24\textwidth}L{0.09\textwidth}@{}}
\caption{Per-task metadata for the 50 MDArena tasks, taken from the \texttt{[metadata]} block of each task's \texttt{task.toml}. Task type and system type are single-valued controlled vocabularies; packages is multi-valued, and a dash means the reference solution uses no domain-specific package.}\label{tab:task-metadata}\\
\toprule
Task ID & Task type & System & Packages & Difficulty \\
\midrule
\endfirsthead
\caption[]{Per-task metadata for the 50 MDArena tasks (continued).}\\
\toprule
Task ID & Task type & System & Packages & Difficulty \\
\midrule
\endhead
\midrule
\multicolumn{5}{r}{Continued on next page}\\
\endfoot
\bottomrule
\endlastfoot
\texttt{01\_\allowbreak{}similes2sim} & System preparation & Small-molecule system (Solvated ligand from SMILES) & GROMACS & Medium \\
\texttt{02\_\allowbreak{}conc\_\allowbreak{}diff\_\allowbreak{}sys\_\allowbreak{}traj} & Diagnostics & Soluble protein (Chignolin; T4 lysozyme) & GROMACS & Easy \\
\texttt{03\_\allowbreak{}id\_\allowbreak{}bp\_\allowbreak{}resid} & Analysis & Membrane protein (Adenosine A2A receptor) & -- & Easy \\
\texttt{04\_\allowbreak{}traj\_\allowbreak{}ana\_\allowbreak{}dist} & Analysis & Membrane protein (Unspecified membrane protein (POPE/POPG bilayer)) & MDAnalysis & Easy \\
\texttt{05\_\allowbreak{}traj\_\allowbreak{}ana\_\allowbreak{}hbonds} & Analysis & Membrane protein (Unspecified membrane protein-ligand complex (POPE/POPG bilayer)) & MDAnalysis & Medium \\
\texttt{06\_\allowbreak{}traj\_\allowbreak{}ana\_\allowbreak{}rotamers} & Analysis & Soluble protein (p38 alpha) & MDAnalysis & Easy \\
\texttt{07\_\allowbreak{}tyk2\_\allowbreak{}equil} & System preparation & Soluble protein (TYK2) & GROMACS & Medium \\
\texttt{08\_\allowbreak{}clash\_\allowbreak{}check} & Diagnostics & Soluble protein (TYK2) & -- & Easy \\
\texttt{09\_\allowbreak{}holo2apo} & System preparation & Membrane protein (mGlu5 receptor) & GROMACS & Medium \\
\texttt{10\_\allowbreak{}traj\_\allowbreak{}ana\_\allowbreak{}mg\_\allowbreak{}coord} & Analysis & Soluble protein (Isocitrate lyase) & MDAnalysis & Easy \\
\texttt{11\_\allowbreak{}amb2gmx\_\allowbreak{}em} & System preparation & Peptide (Glycosylated insulin) & GROMACS & Easy \\
\texttt{12\_\allowbreak{}xtal\_\allowbreak{}waters} & System preparation & Soluble protein (Isocitrate lyase) & MDAnalysis & Medium \\
\texttt{13\_\allowbreak{}altloc\_\allowbreak{}clean} & System preparation & Soluble protein (Oxymyoglobin) & -- & Easy \\
\texttt{14\_\allowbreak{}traj\_\allowbreak{}ana\_\allowbreak{}lig\_\allowbreak{}rmsd} & Analysis & Membrane protein (SLCO2A1) & MDAnalysis & Easy \\
\texttt{15\_\allowbreak{}meng\_\allowbreak{}coupling\_\allowbreak{}groups} & Diagnostics & Membrane protein (MenG) & GROMACS & Medium \\
\texttt{16\_\allowbreak{}gaff\_\allowbreak{}param} & Parameterization & Small-molecule system (Ala-Phe ligand) & ACPYPE; \allowbreak{}GROMACS & Medium \\
\texttt{17\_\allowbreak{}prolif\_\allowbreak{}int} & Analysis & Soluble protein (Isocitrate lyase) & ProLIF & Medium \\
\texttt{18\_\allowbreak{}scale\_\allowbreak{}charge} & Parameterization & Small-molecule system (Water-decane interface with ions) & GROMACS & Easy \\
\texttt{19\_\allowbreak{}traj\_\allowbreak{}ana\_\allowbreak{}open} & Analysis & Membrane protein (TPC2) & MDAnalysis & Medium \\
\texttt{20\_\allowbreak{}fep\_\allowbreak{}abfe\_\allowbreak{}dhdl\_\allowbreak{}analysis} & Free energy & Membrane protein (Adenosine A2A receptor) & alchemlyb; \allowbreak{}GROMACS & Hard \\
\texttt{21\_\allowbreak{}water\_\allowbreak{}transfer\_\allowbreak{}ox2r} & System preparation & Membrane protein (OX2R receptor) & GROMACS; \allowbreak{}MDAnalysis & Hard \\
\texttt{22\_\allowbreak{}traj\_\allowbreak{}ana\_\allowbreak{}pept2\_\allowbreak{}pca\_\allowbreak{}cv\_\allowbreak{}suitability} & Enhanced sampling & Membrane protein (PepT2) & GROMACS; \allowbreak{}PLUMED & Hard \\
\texttt{23\_\allowbreak{}a2ar\_\allowbreak{}membrane\_\allowbreak{}equilibration} & System preparation & Membrane protein (Adenosine A2A receptor) & ACPYPE; \allowbreak{}GROMACS; \allowbreak{}MDAnalysis; \allowbreak{}Open Babel; \allowbreak{}OpenMM; \allowbreak{}PDBFixer & Hard \\
\texttt{24\_\allowbreak{}ligand\_\allowbreak{}parameterization\_\allowbreak{}smirnoff\_\allowbreak{}openfe} & Parameterization & Small-molecule system (Generic ligand for SMIRNOFF parameterization) & GROMACS; \allowbreak{}OpenFF Toolkit; \allowbreak{}OpenMM; \allowbreak{}openmmforcefields; \allowbreak{}ParmEd & Hard \\
\texttt{25\_\allowbreak{}mglu5\_\allowbreak{}membrane\_\allowbreak{}setup} & System preparation & Membrane protein (mGlu5 receptor) & GROMACS; \allowbreak{}OpenFF Toolkit; \allowbreak{}OpenMM; \allowbreak{}ParmEd; \allowbreak{}PDBFixer; \allowbreak{}RDKit & Hard \\
\texttt{26\_\allowbreak{}d3\_\allowbreak{}receptor\_\allowbreak{}membrane\_\allowbreak{}setup} & System preparation & Membrane protein (Dopamine D3 receptor) & ACPYPE; \allowbreak{}GROMACS; \allowbreak{}MDAnalysis; \allowbreak{}Open Babel; \allowbreak{}OpenMM; \allowbreak{}PDBFixer & Hard \\
\texttt{27\_\allowbreak{}a2ar\_\allowbreak{}6gt3\_\allowbreak{}md\_\allowbreak{}setup} & System preparation & Membrane protein (Adenosine A2A receptor) & ACPYPE; \allowbreak{}GROMACS; \allowbreak{}MDAnalysis; \allowbreak{}Open Babel; \allowbreak{}OpenMM; \allowbreak{}PDBFixer & Hard \\
\texttt{28\_\allowbreak{}traj\_\allowbreak{}ana\_\allowbreak{}pept2\_\allowbreak{}pca\_\allowbreak{}plumed\_\allowbreak{}setup} & Enhanced sampling & Membrane protein (PepT2) & GROMACS; \allowbreak{}PLUMED & Hard \\
\texttt{29\_\allowbreak{}pept2\_\allowbreak{}beta\_\allowbreak{}lactam\_\allowbreak{}membrane\_\allowbreak{}setup} & System preparation & Membrane protein (PepT2) & ACPYPE; \allowbreak{}GROMACS; \allowbreak{}Open Babel; \allowbreak{}RDKit & Hard \\
\texttt{30\_\allowbreak{}cam\_\allowbreak{}umbrella\_\allowbreak{}pore\_\allowbreak{}6nq0} & Enhanced sampling & Membrane protein (6NQ0 pore) & GROMACS; \allowbreak{}MDAnalysis; \allowbreak{}PLUMED & Hard \\
\texttt{31\_\allowbreak{}traj\_\allowbreak{}ana\_\allowbreak{}na\_\allowbreak{}conductance\_\allowbreak{}6nq0} & Analysis & Membrane protein (6NQ0 pore) & MDAnalysis & Hard \\
\texttt{32\_\allowbreak{}fep\_\allowbreak{}openfe\_\allowbreak{}rbfe\_\allowbreak{}dhodh} & Free energy & Soluble protein (DHODH) & OpenFE; \allowbreak{}OpenFF Toolkit; \allowbreak{}RDKit & Medium \\
\texttt{33\_\allowbreak{}fep\_\allowbreak{}openfe\_\allowbreak{}rbfe\_\allowbreak{}d3} & Free energy & Membrane protein (Dopamine D3 receptor) & GROMACS; \allowbreak{}OpenFE; \allowbreak{}OpenFF Toolkit; \allowbreak{}OpenMM; \allowbreak{}RDKit & Medium \\
\texttt{34\_\allowbreak{}fep\_\allowbreak{}openfe\_\allowbreak{}rbfe\_\allowbreak{}def\_\allowbreak{}ox2\_\allowbreak{}s1} & Free energy & Membrane protein (OX2R receptor) & GROMACS; \allowbreak{}OpenFE; \allowbreak{}OpenFF Toolkit; \allowbreak{}OpenMM; \allowbreak{}RDKit & Medium \\
\texttt{35\_\allowbreak{}fep\_\allowbreak{}openfe\_\allowbreak{}abfe\_\allowbreak{}d3} & Free energy & Membrane protein (Dopamine D3 receptor) & GROMACS; \allowbreak{}OpenFE; \allowbreak{}OpenFF Toolkit; \allowbreak{}OpenMM; \allowbreak{}RDKit & Medium \\
\texttt{36\_\allowbreak{}fep\_\allowbreak{}openfe\_\allowbreak{}abfe\_\allowbreak{}mglu5r} & Free energy & Membrane protein (mGlu5 receptor) & GROMACS; \allowbreak{}OpenFE; \allowbreak{}OpenFF Toolkit; \allowbreak{}OpenMM; \allowbreak{}RDKit & Medium \\
\texttt{37\_\allowbreak{}fep\_\allowbreak{}openfe\_\allowbreak{}abfe\_\allowbreak{}a2ar} & Free energy & Membrane protein (Adenosine A2A receptor) & GROMACS; \allowbreak{}MDAnalysis; \allowbreak{}OpenFE; \allowbreak{}OpenFF Toolkit; \allowbreak{}OpenMM; \allowbreak{}RDKit & Medium \\
\texttt{38\_\allowbreak{}pept2\_\allowbreak{}qm\_\allowbreak{}cluster\_\allowbreak{}orca} & Quantum chemistry & Membrane protein (PepT2) & MDAnalysis; \allowbreak{}ORCA & Hard \\
\texttt{39\_\allowbreak{}fep\_\allowbreak{}gmx\_\allowbreak{}d3\_\allowbreak{}abfe} & Free energy & Membrane protein (Dopamine D3 receptor) & GROMACS; \allowbreak{}MDAnalysis; \allowbreak{}MDRestraintsGenerator & Medium \\
\texttt{40\_\allowbreak{}traj\_\allowbreak{}ana\_\allowbreak{}meng\_\allowbreak{}sam\_\allowbreak{}rmsd} & Analysis & Membrane protein (MenG) & GROMACS & Medium \\
\texttt{41\_\allowbreak{}rexee\_\allowbreak{}argon\_\allowbreak{}hydration\_\allowbreak{}setup} & Enhanced sampling & Small-molecule system (Argon in water) & ensemble\_md; \allowbreak{}GROMACS & Medium \\
\texttt{42\_\allowbreak{}expanded\_\allowbreak{}ensemble\_\allowbreak{}mdp} & Enhanced sampling & None & GROMACS & Medium \\
\texttt{43\_\allowbreak{}cb7\_\allowbreak{}10\_\allowbreak{}weight\_\allowbreak{}updating\_\allowbreak{}ee} & Enhanced sampling & Small-molecule system (CB7-10 host-guest complex) & GROMACS & Medium \\
\texttt{44\_\allowbreak{}traj\_\allowbreak{}ana\_\allowbreak{}atp8b\_\allowbreak{}rmsd} & Analysis & Membrane protein (ATP8B1-CDC50A complex) & GROMACS; \allowbreak{}MDAnalysis & Medium \\
\texttt{45\_\allowbreak{}atp8b\_\allowbreak{}insertion\_\allowbreak{}depth} & Analysis & Membrane protein (ATP8B1-CDC50A complex) & -- & Hard \\
\texttt{46\_\allowbreak{}meng\_\allowbreak{}dmk\_\allowbreak{}ff\_\allowbreak{}patch} & Parameterization & Membrane protein (MenG) & GROMACS & Hard \\
\texttt{47\_\allowbreak{}d2p\_\allowbreak{}residue\_\allowbreak{}parameterization} & Parameterization & Peptide (Dianionic phosphoaspartate pentapeptide) & GROMACS & Hard \\
\texttt{48\_\allowbreak{}cb7\_\allowbreak{}10\_\allowbreak{}alchemical\_\allowbreak{}metadynamics} & Enhanced sampling & Small-molecule system (CB7-10 host-guest complex) & GROMACS; \allowbreak{}PLUMED & Medium \\
\texttt{49\_\allowbreak{}slco2a1\_\allowbreak{}physical\_\allowbreak{}validation} & Diagnostics & Membrane protein (SLCO2A1) & GROMACS; \allowbreak{}physical\_validation & Medium \\
\texttt{50\_\allowbreak{}glycoinsulin\_\allowbreak{}beta\_\allowbreak{}sheet\_\allowbreak{}propensity} & Analysis & Peptide (Glycosylated insulin) & GROMACS; \allowbreak{}MDAnalysis & Medium \\
\end{longtable}
\endgroup
\clearpage


\begin{table}[!ht]
\centering
\caption{Positioning of MDArena relative to representative agentic molecular-simulation and computational-chemistry systems. \cmark, \pmark, and \xmark{} denote yes, partial or limited coverage, and no, respectively.}
\label{tab:benchmark-comparison}
\footnotesize
\setlength{\tabcolsep}{5pt}
\begin{tabular}{@{}L{0.36\linewidth}cccc@{}}
\toprule
System or benchmark & \shortstack{Biomol.\\MD} & \shortstack{Exec.\\workflows} & \shortstack{Multi-\\package} & \shortstack{Research\\provenance} \\
\midrule
MDCrow~\citep{campbell2026mdcrow} & \cmark & \cmark & \pmark & \pmark \\
NAMD-Agent~\citep{chandrasekhar2025namdagent} & \cmark & \cmark & \xmark & \pmark \\
DynaMate~\citep{guilbert2025dynamate} & \cmark & \cmark & \pmark & \pmark \\
PolyJarvis~\citep{zhao2026polyjarvis} & \xmark & \cmark & \xmark & \pmark \\
LAMMPS-Agents~\citep{vriza2026multiagentic} & \xmark & \cmark & \xmark & \pmark \\
ChemGraph~\citep{pham2026chemgraph} & \pmark & \pmark & \cmark & \pmark \\
MD-EvalBench~\citep{shi2026mdagent2} & \xmark & \pmark & \xmark & \pmark \\
MDGym~\citep{kumar2026mdgym} & \pmark & \cmark & \pmark & \pmark \\
\textbf{MDArena} & \cmark & \cmark & \cmark & \cmark \\
\bottomrule
\end{tabular}
\end{table}



\clearpage
\section{Supplementary Figures}
\begin{figure}[!ht]
\centering
\suppfigure{figures/figure_S1/figure_S1.pdf}
\caption{Performance and cost of the nine coding-agent configurations omitted from main-text Figure~3, using the same analysis and plotting style. (A) Mean Strict-Pass@1 by task difficulty. (B) Mean cost per trial by difficulty. (C) Overall Strict-Pass@1 against mean cost per trial. A strict pass requires a completed trial with correctness and process scores of one; errors count as failures. Error bars show 95\% confidence intervals from 2000 task-bootstrap resamples. Superscript stars ($\ast$) mark MDArena v0.2 configurations with three attempts per task; daggers ($\dagger$) mark historical v0.1 configurations with one attempt per task. The v0.1 results retain their original task definitions and difficulty labels and are historical context rather than a controlled comparison across benchmark versions. Recorded costs are averaged over available cost records. Striped Qwen3 Coder 480B A35B bars use the workshop evaluation's supplied mean cost of GBP~0.14 per trial, converted at 1~GBP~=~1.36~USD to USD~0.1904. Its costs by difficulty are estimated from token counts using relative uncached-input, cached-input, and output prices of 1:0.1:4, scaled to preserve that supplied overall mean.}
\label{fig:supp-difficulty-cost}
\end{figure}

The following figures report the seven-configuration workshop evaluation,
with three attempts per task and 150 trials per configuration. They do not
represent the newer model evaluations reported in the current main text.
% Replace these figures and captions together when updated assets are available.

\begin{figure}[!ht]
\centering
\suppfigure{figures/codex_regrade_human/figure_A1.pdf}
\caption{Trial outcomes in the workshop evaluation by model, as a percentage of each configuration's 150 trials. A crashed trial is counted as an error even when the verifier still scored it, so the full-pass band equals Strict-Pass@1.}
\label{fig:outcome-stacked-by-model}
\end{figure}

\begin{figure}[!ht]
\centering
\suppfigure{figures/codex_regrade_human/figure_A2.pdf}
\caption{Performance by task category in the workshop evaluation: Strict-Pass@1 (left), correctness (middle) and process (right) scores, macro-averaged over tasks with errored trials scored zero. The number of tasks per category is given in the row label. Error bars are omitted: several categories hold too few tasks to support one.}
% Both labels point here: this figure absorbed the main-text category figure, and
% the Results prose \cref{}s each of them.
\label{fig:category-check-scores}
\label{fig:success-by-task-category}
\end{figure}


\clearpage
% Use the same RSC bibliography style and source as the main manuscript.
\IfFileExists{rsc.bib}{%
  \bibliographystyle{rsc}
  \bibliography{rsc}
}{%
  \section*{References}
  \addcontentsline{toc}{section}{References}
  \noindent\textcolor{red!65!black}{The shared bibliography file \texttt{rsc.bib} was not found.}
}
\end{document}
